\subsection{主元素}

定义 3。假设 g 是半单的。

\begin{enumerate}
\def\labelenumi{(\roman{enumi})}
\item
  若 \(x\) 的幂零元素 \(\mathfrak{g}\) 的 \(\operatorname{Ker} \operatorname{ad} x\) 的维数等于 \(\mathfrak{g}\) 的秩，则称 x 为主幂零元素。
\item
  若 \(h\)\(\mathfrak{g}\) 的单元素 \(h\) 是正则的，且 \(\operatorname{ad} h\) 在 k 的代数闭包中的特征值属于 2Z，则称 h 为主单元素。\(k\)
\item
  若  的 \(\mathfrak{sl}_2\) 三元组 \((x,h,y)\) 在视为 \(\mathfrak{g}\)\(\mathfrak{g}\)\(kx + kh + ky\) 上的模时的长度等于 \(\mathfrak{g}\) 的秩，则称其为主三元组。
\end{enumerate}

命题 7。假设 g 是半单的。令 \((x,h,y)\) 为 g 中的一个 \(sl_{2}\) 三元组。以下条件等价：

\begin{enumerate}
\def\labelenumi{(\roman{enumi})}
\item
  \(x\) 是主幂零元素；
\item
  h 是主单元素；
\item
  \((x,h,y)\) 是主三元组。
\end{enumerate}

对 \(p \in Z\)，令 \(\mathfrak{g}^{p} = \operatorname{Ker}(\operatorname{ad} h - p)\)。令 \(g' = \sum_{p \in Z} g^{2p}\)。若将 g 视为 \(a = kx + kh + ky\) 上的模，则 \(g'\) 是奇数维单子模之和（§1，第2号，命题2的推论）。令 l（分别为 \(l'\)）为将 g（分别为 \(g'\)）视为 a-模时的长度。由 §1，第2号，

\[
\dim (\operatorname{Ker} \operatorname{ad} x) = l \geq l ^ {\prime} = \dim (\operatorname{Ker} \operatorname{ad} h) \geq \operatorname{rk} (\mathfrak {g}).
\]

(i) 与 (iii) 的等价性立即得到。另一方面，条件 (ii) 意味着 \(\dim(\operatorname{Ker} \operatorname{ad} h) = \operatorname{rk}(\mathfrak{g})\) 且 \(g' = g\)，换言之，

\[
\dim (\operatorname{Ker} \operatorname{ad} h) = \operatorname{rk} (\mathfrak {g})
\]

且 \(l = l'\)。条件 (ii) 与其他条件的等价性随即得到。

命题 8。假设 \(\neq 0\) 是非零半单李代数。令  为 g 的分裂嘉当子代数，R 为 \((\mathfrak{g},\mathfrak{h})\) 的根系，B 为 R 的一个基，令 \(h^{0}\) 为 h 中满足对所有  都有 \(\alpha(h^{0}) = 2\) 的元素。\(\alpha \in B\)

\begin{enumerate}
\def\labelenumi{(\roman{enumi})}
\item
  元素 \(h^0\) 是单元素且是主单元素。
\item
  \(x\)\(\mathfrak{g}\) 中满足存在形如  的 \(\mathfrak{sl}_2\) 三元组的元素 x，正是 \((x, h^0, y)\)\(\sum_{\alpha \in B} \mathfrak{g}^\alpha\) 中在每个 \(\mathfrak{g}^\alpha\) 中都有非零分量的元素。
\end{enumerate}

元素 \(h^0\) 就是 §7 第5号引理2 中考察的元素（参见同处公式 (1)）。由该引理可知，\(h^0\) 是主单元素，并且若 \(x \in \sum_{\alpha \in \mathrm{B}} \mathfrak{g}^\alpha\) 在每个 \(\mathfrak{g}^\alpha\) 中都有非零分量，则存在形如 \(\mathfrak{sl}_2\) 的 \((x, h^0, y)\) 三元组。反之，令 \((x, h^0, y)\) 为一个 \(\mathfrak{sl}_2\) 三元组。有 \([h^0, x] = 2x\)，所以 \(x \in \sum_{\gamma \in \mathrm{R}, \gamma(h^0) = 2} \mathfrak{g}^\gamma = \sum_{\alpha \in \mathrm{B}} \mathfrak{g}^\alpha\)。同样，\(y \in \sum_{\alpha \in \mathrm{B}} \mathfrak{g}^{-\alpha}\)。写

\[
h ^ {0} = \sum_ {\alpha \in \mathrm{B}} a _ {\alpha} H _ {\alpha} \quad \text { where } a _ {\alpha} > 0 \text { for all } \alpha \in \mathrm{B},
\]

\[
x = \sum_ {\alpha \in \mathrm{B}} X _ {\alpha} \qquad \text { where } X _ {\alpha} \in \mathfrak {g} ^ {\alpha} \text { for   all } \alpha \in \mathrm{B},
\]

\[
y = \sum_ {\alpha \in \mathrm{B}} X _ {- \alpha} \quad \text { where } X _ {- \alpha} \in \mathfrak {g} ^ {- \alpha} \text { for   all } \alpha \in \mathrm{B}.
\]

则

\[
\sum_ {\alpha \in \mathrm{B}} a _ {\alpha} H _ {\alpha} = h ^ {0} = [ y, x ] = \sum_ {\alpha , \beta \in \mathrm{B}} [ X _ {- \beta}, X _ {\alpha} ] = \sum_ {\alpha \in \mathrm{B}} [ X _ {- \alpha}, X _ {\alpha} ]
\]

所以对所有 \([X_{-\alpha}, X_{\alpha}] \neq 0\) 都有 \(\alpha \in \mathrm{B}\)。

推论。在可分裂半单李代数中，存在主幂零元素。

在非可分裂半单李代数中，0 可能是唯一的幂零元素。

命题 9。假设 \(k\) 是代数闭域且 \(\mathfrak{g}\) 是半单的。所有 \(\mathfrak{g}\) 的主单元素（分别为主幂零元素）都在 \(\operatorname{Aut}_e(\mathfrak{g})\) 下共轭。

沿用命题8的记号。令 h 为一个主单元素。它在 \(\operatorname{Aut}_{e}(\mathfrak{g})\) 下与 h 中的某个 \(h' \in h\) 共轭，使得对所有 \(\alpha(h') \in \{0,1,2\}\) 都有 \(\alpha \in B\)（第3号，命题5）。由于 \(h'\) 是主单元素，对所有 \(\alpha(h') \neq 0\) 都有 \(\alpha(h') \in 2Z\) 且 \(\alpha \in B\)，所以对所有 \(\alpha(h') = 2\) 都有 \(\alpha \in B\)，从而 \(h' = h^{0}\)。这证明了主单元素的断言。

令 \(x, x'\) 为主幂零元素。存在 \(\mathfrak{sl}_2\) 三元组 \((x, h, y)\)、\((x', h', y')\)。由命题7，\(h\) 和 \(h'\) 是主单元素，因而由上文它们在 \(\mathrm{Aut}_e(\mathfrak{g})\) 下共轭。所以 \(x\) 和 \(x'\) 在 \(\mathrm{Aut}_e(\mathfrak{g})\) 下共轭（命题6）。

引理 9。沿用命题8的记号。对 \(\mathfrak{g}^p = \mathrm{Ker}(\mathrm{ad}h^0 - p)\)，置 \(p \in \mathbf{Z}\)。令 \(\mathfrak{g}_*^2\) 为 \(\mathfrak{g}^2 = \sum_{\alpha \in \mathrm{B}} \mathfrak{g}^\alpha\) 中在每个 \(\mathfrak{g}^\alpha\) 中都有非零分量的元素所成的集合。令 \(\mathrm{R}_+\) 为关于 \(\mathrm{B}\) 的正根集合，\(\mathfrak{n}_+ = \sum_{\alpha \in \mathrm{R}_+} \mathfrak{g}^\alpha\)，并令 \(x \in \mathfrak{g}_*^2\)。则 \(e^{\mathrm{ad}\mathfrak{n}_+}. x = x + [\mathfrak{n}_+, \mathfrak{n}_+]\)。

显然 \(e^{\mathrm{ad} \mathfrak{n}_+}.x \subset x + [\mathfrak{n}_+, \mathfrak{n}_+]\)。我们证明，若 \(v \in [\mathfrak{n}_+, \mathfrak{n}_+]\)，则 \(x + v \in e^{\mathrm{ad} \mathfrak{n}_+}.x\)。置 \(\mathfrak{n}^{(p)} = \sum_{r \geq p} \mathfrak{g}^{2r}\)；只需证明

\[
x + v \in e ^ {\mathrm{ad} \mathfrak {n} _ {+}}. x + \mathfrak {n} ^ {(p)}
\]

对所有 \(p \geq 2\) 都成立。当 \(p = 2\) 时这显然成立，因为 \(\mathfrak{n}^{(2)} = [\mathfrak{n}_+, \mathfrak{n}_+]\)（§3，第3号，命题9 (iii)）。假设我们已经找到 \(z \in \mathfrak{n}_+\)，使 \(v + x - e^{\mathrm{ad}z}.x \in \mathfrak{n}^{(p)}\)。由于存在形如 \(\mathfrak{sl}_2\) 的 \((x, h^0, y)\) 三元组（命题8），§1 第2号命题2的推论说明 \([x, \mathfrak{g}^{2p-2}] = \mathfrak{g}^{2p}\)；因此存在 \(z' \in \mathfrak{g}^{2p-2} \subset \mathfrak{n}_+\)，使

\[
v + x - e ^ {\mathrm{ad} z}. x \in [ z ^ {\prime}, x ] + \mathfrak {n} ^ {(p + 1)}.
\]

所以 \(v + x \in e^{\mathrm{ad}(z + z')}.x + \mathfrak{n}^{(p + 1)}\)，归纳即证明了断言。

命题 10。假设 g 是半单的。令  为 g 的分裂嘉当子代数，R 为 \((\mathfrak{g},\mathfrak{h})\) 的根系，B 为 R 的一个基，\(R_{+}\) 为关于 B 的正根集合，且 \(n_{+}=\sum_{\alpha\in R_{+}}g^{\alpha}\)。属于 \(n_{+}\) 的主幂零元素正是 \(n_{+}\) 中在每个  对应的 \(g^{\alpha}\) 中都有非零分量的元素。\(\alpha\in B\)

命题8和引理9证明了这类元素是主幂零元素。我们证明逆命题。显然可以假设 \(\mathfrak{g}\) 是单的。令 \(h^0\) 和 \(\mathfrak{g}^p\) 如命题8和引理9中所述。令 \(\omega\) 为最高根，并置 \(\omega(h^0) = 2q\)；有 \(q = h - 1\)，其中 h 为 R 的 Coxeter 数，参见~第 VI 章第1节第11号命题31。于是 \(h\)\(\mathfrak{g}^{2q} = \mathfrak{g}^\omega\)、\(\mathfrak{g}^{-2q} = \mathfrak{g}^{-\omega}\)，且当  时 \(\mathfrak{g}^{2k} = 0\)。存在一个主 \(|k| > q\)\(\mathfrak{sl}_2\) 三元组 \((x^0, h^0, y^0)\)\((\mathrm{ad} x^0)^{2q}(\mathfrak{g}^{-\omega}) = \mathfrak{g}^\omega\)\((\mathrm{ad} x^0)^{2q} \neq 0\)\(x\)。由 §1 第2号命题2的推论，，所以 。令 x 为属于  的 \(\mathfrak{g}\) 的主幂零元素。若 \(\mathfrak{n}_+\)\(\bar{k}\) 是 k 的代数闭包，则 \(k\)\(x \otimes 1\) 和 \(x^0 \otimes 1\) 在 \(\mathfrak{g} \otimes_k \bar{k}\) 的某个自同构下共轭（命题9），所以 \((\mathrm{ad} x)^{2q} \neq 0\)。存在 \(\lambda \in \mathbb{R}\)，使 \((\mathrm{ad} x)^{2q} \mathfrak{g}^\lambda \neq 0\)。置 \(x = \sum_{n \geq 1} x_n\)，其中 \(x_n \in \mathfrak{g}^{2n}\)。则

\[
(\operatorname{ad} x) ^ {2 q} \mathfrak {g} ^ {\lambda} \subset (\operatorname{ad} x _ {1}) ^ {2 q} \mathfrak {g} ^ {\lambda} + \sum_ {k > 4 q + \lambda (h ^ {0})} \mathfrak {g} ^ {k} = (\operatorname{ad} x _ {1}) ^ {2 q} \mathfrak {g} ^ {\lambda},
\]

由于 \(4q + \lambda(h^0) \geq 4q - 2q = 2q\)。现在，当 \((\operatorname{ad} x_1)^{2q} \mathfrak{g}^\lambda \subset \mathfrak{g}^{4q + \lambda(h^0)}\) 时，\(\lambda = -\omega\)。因此，\((\operatorname{ad} x_1)^{2q} \mathfrak{g}^{-\omega} = \mathfrak{g}^\omega\)。我们有 \(\omega = \sum_{\alpha \in \mathrm{B}} n_\alpha \alpha\)，其中对所有 \(n_\alpha > 0\) 都有 \(\alpha \in \mathrm{B}\)（第 VI 章第1节第8号注）。若存在 \(\alpha_0 \in \mathrm{B}\)，使 \(x_1 \in \sum_{\alpha \in \mathrm{B}, \alpha \neq \alpha_0} \mathfrak{g}^\alpha\)，则关系

\[
\omega \notin - \omega + \sum_ {\alpha \in \mathrm{B}, \alpha \neq \alpha_ {0}} k \alpha
\]

这意味着对所有 p 都有 \(\mathfrak{g}^{\omega} \not\subset (\operatorname{ad} x_1)^p \mathfrak{g}^{-\omega}\)\(p\)\(x_1\)，这是荒谬的。因此，x₁ 在每个 \(\mathfrak{g}^{\alpha}\) 中的分量对所有 \(\alpha \in B\) 都非零。

